\documentclass[12pt]{article}
\usepackage[utf8]{inputenc}
\usepackage[T1]{fontenc}
\usepackage{hyperref}
\title{Tsamai Language Resources for K-12 Curriculum Development}
\author{Compiled by Claude Sonnet 5 (claude-sonnet-5) for Andr\'as Kornai}
\date{}
\begin{document}
\maketitle

\section{Introduction}
Tsamai (Ts'amakko; ISO 639-3 tsb; Glottocode tsam1247) is a Dullay Cushitic language of southwestern Ethiopia. Its speakers live mostly in Bena Tsemay woreda, Debub Omo Zone. Speaker counts conflict: 18,000 (2007 census, per Wikipedia), 8,621 (1998 census, per the Tsamai people article) and 19,200 (Ethnologue, per German Wikipedia). Search results found little educational material.

\section{Primary Educational Resources}
No primers, textbooks or school programmes were found.

\section{Linguistic and Reference Materials}
\begin{itemize}
\item Tsamai language (Wikipedia). It cites Graziano Sav\`a, \emph{A Grammar of Ts'amakko} (Cushitic Language Studies 22, 2005), the main reference grammar: \url{https://en.wikipedia.org/wiki/Tsamai_language}
\item Tsamai people (Wikipedia): \url{https://en.wikipedia.org/wiki/Tsamai_people}
\item Tsamay (German Wikipedia): \url{https://de.wikipedia.org/wiki/Tsamay_(Sprache)}
\end{itemize}
The grammar itself was not located online in these searches.

\section{Religious and Community Resources}
None found.

\section{Digital and Audio Resources}
\begin{itemize}
\item Wikitongues page for Tsamai: \url{https://wikitongues.org/languages/tsb/}
\item Endangered Languages Project entry: \url{https://endangeredlanguages.com/node/38443}
\item Verbix entry: \url{https://docs.verbix.com/EndangeredLanguages/Tsamay}
\end{itemize}

\section{Curriculum Development Context}
Sources describe the language as largely unwritten. German Wikipedia says it uses the Latin script. A curriculum would build on Sav\`a's grammar and new community-developed materials.
\end{document}